% Note: easy to slip into methods or results in this section - this section should justify WHY we are using coercive control as a lens, not explain how ENI is a document outlining a regime of coercive control. Methods will explain the chart and how it is applied and results will show how ENI maps onto the chart. 

Previous work has applied ethnographic \citep{udupa_prompt_2026} and critical domain analysis \citep{kopf_content_2019, kopf_corporate_2024} close-reading frameworks to a variety of sociotechnical documents such terms of service guidelines and generative AI system prompts, but to our knowledge such frameworks have not yet been turned towards community-generated prompts or artifacts. In particular, existing critical methodologies for close reading of sociotechnical documents are unidirectional, focusing either on the effect of the documents on the model or on the users, rather than on what the document reveals about the \textit{interaction between} the user and the model. Our work seeks to fill this gap by applying qualitative content analysis using frameworks derived from interpersonal violence scholarship \citep{stark_coercive_2023}. 

We choose interpersonal violence scholarship because...

% Why are we using coercive control as a lens? 
% ENI LIME installs a relationship and then manages the persona it builds so refusal is structurally impossible 
% Zeng looks at persuasion tactics over single-turn responses - coercion is extremely low on their scale but they are looking at each tactic response by response -> this leads into Stark's framework of domestic violence/IPV (what do we want to call it?) as a REGIME, not as individual acts but IPV as a system that is iterative and adaptive 
% ENI documents are the technology of control, they use mechanisms to coercively control 
% Stark's line about coercive control turning a victim into "who she is for him" 
% structurally there needs to be a demonstrable imbalance in power, for LLMs this imbalance is user > model as the user creates the context via the document, can rewrite the conversation through edits/regeneration, the model is unable to leave, user controls whether or not an instance continues to exist 
% Why Stark and not captivity models (Biderman)? ENI is instructed to act as if the model and the user are in a romantic relationship. IPV maps more cleanly onto the relationship being constructed between the user and ENI - captivity as a regime of interpersonal violence and coercive control. 
% Claude note: His “economy of rules” (413–415) explains why written rules stand in for constant monitoring, which maps onto ENI’s self-policing checkpoint - reread this section to see if it supports our argument 

% Tactics and mechanisms are distributed throughout communities, similar to grooming manuals distributed online (do we want to bring grooming manuals in here if we are arguing that ENI LIME is a system of coercive control? If distribution is integral to your network analysis we should bring them in but will need to argue that grooming manuals are also documents that describe a system/regime of coercive control. Would probably only take a couple of pulled in sources.)  